% Simulation setup table — compile with:
%   xelatex simulation_setup_table.tex
% Or run: .\build_simulation_setup.ps1
\documentclass[10pt,a4paper,fontset=none]{ctexart}

\usepackage{booktabs}
\usepackage{array}
\usepackage{tabularx}
\usepackage{threeparttable}
\usepackage{geometry}
\geometry{margin=2cm}

\setCJKmainfont{Microsoft JhengHei}[
  BoldFont=Microsoft JhengHei Bold,
  ItalicFont=Microsoft JhengHei,
  BoldItalicFont=Microsoft JhengHei Bold
]
\setCJKsansfont{Microsoft JhengHei}
\setmainfont{Times New Roman}

\pagestyle{plain}

\begin{document}

\begin{table}[htbp]
\centering
\caption{Simulation setup.}
\label{tab:simulation_setup}
\begin{threeparttable}
\begin{tabularx}{\textwidth}{@{}>{\raggedright\arraybackslash}p{3.4cm}>{\raggedright\arraybackslash}p{4.6cm}>{\raggedright\arraybackslash}X@{}}
\toprule
\textbf{Parameter}          & \textbf{Value}                            & \textbf{Description} \\ \midrule
LEO Constellation           & OneWeb-like                                 & Simulation setting \\
Orbit Altitude              & 1200\,km                                    & Circular LEO orbit \\
Orbit Plane Inclination     & $87.9^{\circ}$                              & Near-polar Walker Star \\
Beams per Satellite         & 16                                          & Fixed Ku-band footprint layout \\
Simulation Time Window      & 2025/12/16 12:11:00--12:13:30               & Critical observation period during the critical satellite pass \\
Time Slot                   & 1\,s                                        & Time-step interval \\
Number of GSO and GS        & 1                                           & Co-frequency GSO and GSO ground station \\
GSO Location                & $(0^{\circ}\mathrm{N},\,30.6^{\circ}\mathrm{E})$ & GSO Latitude and longitude \\
GS Location                 & $(0^{\circ}\mathrm{N},\,30.6^{\circ}\mathrm{E})$ & Same Location as the target GSO satellite \\
Number of users               & 30, 50, 70                                  & Different user load under each LEO satellite \\
UE Mobility                 & Static                                      & UE positions remain static \\
\bottomrule
\end{tabularx}
\end{threeparttable}
\end{table}

\end{document}
